


\documentclass{ximera}

\input{../../../preamble.tex}


\definecolor{fvdnavy}{HTML}{071D43}
\definecolor{fvdcyan}{HTML}{20BFC6}
\definecolor{fvdcyanlight}{HTML}{EFFCFB}
\definecolor{fvdmagenta}{HTML}{F10D69}
\definecolor{fvdmagentalight}{HTML}{FFF1F7}
\definecolor{fvdtext}{HTML}{163044}





%=================================================
% Pedagogische omgevingen
% PDF: tcolorbox
% HTML: learning-box
%=================================================

\newenvironment{voorkennis}
{%
    \ifdefined\HCode
        \HCode{%
<section class="learning-box learning-box--prior">
<div class="learning-box__header">
<span class="learning-box__icon" aria-hidden="true"></span>
<span class="learning-box__title">Voorkennis</span>
</div>
<div class="learning-box__content">
        }%
        \begin{itemize}
    \else
       \begin{tcolorbox}[
    enhanced,
    breakable,
    colback=fvdcyanlight,
    colframe=fvdcyan,
    coltext=fvdtext,
    boxrule=0.5pt,
    leftrule=4pt,
    arc=4mm,
    outer arc=4mm,
    left=7mm,
    right=6mm,
    top=5mm,
    bottom=5mm,
    before skip=6mm,
    after skip=6mm,
    title=Voorkennis,
    coltitle=fvdcyan,
    fonttitle=\bfseries\Large,
    attach title to upper={\par\smallskip},
    boxed title style={
        empty,
        left=0mm,
        right=0mm,
        top=0mm,
        bottom=0mm
    },
    overlay={
        \draw[
            fvdcyan,
            line width=0.5pt
        ]
        ([xshift=42mm,yshift=-13mm]frame.north west)
        --
        ([xshift=-7mm,yshift=-13mm]frame.north east);
    }
]
        \begin{itemize}
    \fi
}
{%
    \end{itemize}
    \ifdefined\HCode
        \HCode{%
</div>
</section>
        }%
    \else
        \end{tcolorbox}
    \fi
}


\newenvironment{leerdoelen}
{%
    \ifdefined\HCode
        \HCode{%
<section class="learning-box learning-box--goals">
<div class="learning-box__header">
<span class="learning-box__icon" aria-hidden="true"></span>
<span class="learning-box__title">Leerdoelen</span>
</div>
<div class="learning-box__content">
        }%
        \begin{itemize}
    \else
\begin{tcolorbox}[
    enhanced,
    breakable,
    colback=fvdmagentalight,
    colframe=fvdmagenta,
    coltext=fvdtext,
    boxrule=0.5pt,
    leftrule=4pt,
    arc=4mm,
    outer arc=4mm,
    left=7mm,
    right=6mm,
    top=5mm,
    bottom=5mm,
    before skip=6mm,
    after skip=6mm,
    title=Leerdoelen,
    coltitle=fvdmagenta,
    fonttitle=\bfseries\Large,
    attach title to upper={\par\smallskip},
    boxed title style={
        empty,
        left=0mm,
        right=0mm,
        top=0mm,
        bottom=0mm
    },
    overlay={
        \draw[
            fvdmagenta,
            line width=0.5pt
        ]
        ([xshift=42mm,yshift=-13mm]frame.north west)
        --
        ([xshift=-7mm,yshift=-13mm]frame.north east);
    }
]
        \begin{itemize}
    \fi
}
{%
    \end{itemize}
    \ifdefined\HCode
        \HCode{%
</div>
</section>
        }%
    \else
        \end{tcolorbox}
    \fi
}



\addPrintStyle{../../..}

\title{Implicatie en equivalentie}

\author{Ingmar Herreman}

\begin{abstract}

Veel wiskundige uitspraken hebben de vorm
\textbf{``als \ldots, dan \ldots''}.

Als een geheel getal deelbaar is door $4$, dan is het even.
Als een driehoek gelijkzijdig is, dan is hij gelijkbenig.
Als $x>3$, dan is $x^2>9$.

Zulke uitspraken noemen we \textbf{implicaties}.

In dit hoofdstuk onderzoeken we zorgvuldig wat een implicatie betekent.
We leren wanneer een implicatie waar of onwaar is en ontdekken waarom
een implicatie iets anders zegt dan haar omgekeerde.

Daarbij maken we kennis met nodige en voldoende voorwaarden en met de
contrapositieve uitspraak. Ten slotte onderzoeken we wanneer twee
voorwaarden precies even sterk zijn en dus een equivalentie vormen.

Deze begrippen zijn fundamenteel voor wiskundig redeneren.
Later zullen we ze gebruiken om redeneringen te analyseren en
wiskundige uitspraken te bewijzen.

\end{abstract}

\begin{document}

% FVD_CHAPTER_DATA_START
% Automatisch gegenereerd uit index.tex.
% Niet handmatig aanpassen.
\fvdchapterdata{1}{Logica — De taal van correct redeneren}{3}
% FVD_CHAPTER_DATA_END



\maketitle


% ============================================================
% VOORKENNIS
% ============================================================

\begin{voorkennis}

  \item Je weet wat een wiskundige uitspraak is.

  \item Je kan de waarheidswaarde van een uitspraak bepalen.

  \item Je kan uitspraken benoemen met letters zoals $p$ en $q$.

  \item Je kent de logische operatoren
        $\neg$, $\land$ en $\lor$.

  \item Je kan eenvoudige samengestelde uitspraken analyseren.

  \item Je kan een waarheidstabel lezen en opstellen.

  \item Je begrijpt dat één tegenvoorbeeld volstaat om een algemene
        bewering te weerleggen.

\end{voorkennis}


% ============================================================
% LEERDOELEN
% ============================================================

\begin{leerdoelen}

  \item Je kan een implicatie herkennen en noteren als $p\Rightarrow q$.

  \item Je kan in een implicatie de voorwaarde en de conclusie herkennen.

  \item Je kan de waarheidswaarde van een implicatie bepalen.

  \item Je kan de waarheidstabel van een implicatie interpreteren.

  \item Je kan uitleggen wanneer een implicatie onwaar is.

  \item Je kan nodige en voldoende voorwaarden herkennen en formuleren.

  \item Je kan bij een implicatie de omgekeerde uitspraak formuleren.

  \item Je kan bij een implicatie de contrapositieve uitspraak formuleren.

  \item Je kan het verschil uitleggen tussen een implicatie en haar omgekeerde.

  \item Je kan onderzoeken of een implicatie en haar omgekeerde waar zijn.

  \item Je kan een equivalentie herkennen en noteren als $p\Leftrightarrow q$.

  \item Je kan equivalenties formuleren met ``als en slechts als''.

  \item Je kan logische uitspraken in nieuwe situaties analyseren en
        je besluit verantwoorden.

\end{leerdoelen}


% ============================================================
\FVDsection{Als \ldots, dan \ldots}
% ============================================================

Beschouw de uitspraak

\begin{center}
Als een geheel getal deelbaar is door $4$, dan is het even.
\end{center}

Deze uitspraak bestaat uit twee delen.

We noemen

\[
p:\quad n\text{ is deelbaar door }4
\]

en

\[
q:\quad n\text{ is even}.
\]

De oorspronkelijke uitspraak zegt dan:

\begin{center}
als $p$, dan $q$.
\end{center}

Symbolisch schrijven we

\[
p\Rightarrow q.
\]


\begin{definitie}

Een \textbf{implicatie} is een samengestelde uitspraak van de vorm

\[
p\Rightarrow q.
\]

We lezen dit als

\begin{center}
``als $p$, dan $q$''
\end{center}

of

\begin{center}
``$p$ impliceert $q$''.
\end{center}

De uitspraak $p$ noemen we de \textbf{voorwaarde}.

De uitspraak $q$ noemen we de \textbf{conclusie}.

\end{definitie}


\begin{voorbeeld}[Voorwaarde en conclusie]

Beschouw

\begin{center}
Als $n$ een veelvoud van $6$ is, dan is $n$ deelbaar door $3$.
\end{center}

We kunnen schrijven

\[
p:\quad n\text{ is een veelvoud van }6
\]

en

\[
q:\quad n\text{ is deelbaar door }3.
\]

De uitspraak heeft dus de vorm

\[
p\Rightarrow q.
\]

De voorwaarde is dat $n$ een veelvoud van $6$ is.

De conclusie is dat $n$ deelbaar is door $3$.

\end{voorbeeld}


\begin{opmerking}

De pijl

\[
\Rightarrow
\]

betekent niet dat we een berekening uitvoeren.

Hij drukt een \textbf{logisch verband} uit:

wanneer de voorwaarde waar is, moet ook de conclusie waar zijn.

\end{opmerking}


\begin{oefening}[Voorwaarde en conclusie herkennen]{\oefB\ESS}

Beschouw de uitspraak:

\begin{center}
Als een getal deelbaar is door $10$, dan is het even.
\end{center}

\begin{question}

Welke uitspraak is de voorwaarde?

\begin{multipleChoice}
  \choice[correct]{Het getal is deelbaar door $10$.}
  \choice{Het getal is even.}
\end{multipleChoice}

\end{question}

\begin{question}

Welke uitspraak is de conclusie?

\begin{multipleChoice}
  \choice{Het getal is deelbaar door $10$.}
  \choice[correct]{Het getal is even.}
\end{multipleChoice}

\end{question}

\begin{question}

Schrijf de uitspraak symbolisch wanneer

\[
p:\quad n\text{ is deelbaar door }10
\]

en

\[
q:\quad n\text{ is even}.
\]

\begin{oplossing}

\[
p\Rightarrow q.
\]

\end{oplossing}

\end{question}

\end{oefening}


% ============================================================
\FVDsection{Wanneer is een implicatie waar?}
% ============================================================

Bij een conjunctie of disjunctie konden we de waarheidswaarde bepalen
uit de waarheidswaarden van $p$ en $q$.

Dat kan ook bij een implicatie.

Beschouw opnieuw

\[
p\Rightarrow q.
\]

De belangrijkste vraag is:

\begin{center}
Wanneer wordt de belofte ``als $p$, dan $q$'' geschonden?
\end{center}


% ------------------------------------------------------------
\FVDsubsection{De voorwaarde is waar}
% ------------------------------------------------------------

Stel dat $p$ waar is.

Dan beweert de implicatie dat ook $q$ waar moet zijn.

Als $q$ inderdaad waar is, klopt de implicatie.

Als $q$ onwaar is, hebben we precies een situatie gevonden waarin
de voorwaarde vervuld is maar de conclusie niet.

Dan is de implicatie onwaar.


\begin{voorbeeld}

Beschouw voor een geheel getal $n$:

\[
p:\quad n\text{ is deelbaar door }4
\]

en

\[
q:\quad n\text{ is even}.
\]

Voor $n=12$ zijn $p$ en $q$ beide waar.

De implicatie

\[
p\Rightarrow q
\]

wordt in dit geval dus niet geschonden.

\end{voorbeeld}


\begin{voorbeeld}[Hoe maak je een implicatie onwaar?]

Beschouw de bewering

\begin{center}
Als een geheel getal even is, dan is het deelbaar door $4$.
\end{center}

Neem

\[
n=6.
\]

Dan is de voorwaarde

\[
6\text{ is even}
\]

waar, maar de conclusie

\[
6\text{ is deelbaar door }4
\]

onwaar.

We hebben dus een geval gevonden waarin

\[
p\text{ waar}
\qquad\text{en}\qquad
q\text{ onwaar}.
\]

Daarom is de algemene implicatie onwaar.

Het getal $6$ is een \textbf{tegenvoorbeeld}.

\end{voorbeeld}


\begin{opmerking}

Om een algemene implicatie

\[
p\Rightarrow q
\]

te weerleggen, zoeken we dus naar een geval waarvoor

\[
p\text{ waar}
\qquad\text{en}\qquad
q\text{ onwaar}.
\]

Eén zo'n tegenvoorbeeld volstaat.

\end{opmerking}


% ------------------------------------------------------------
\FVDsubsection{De voorwaarde is onwaar}
% ------------------------------------------------------------

Wat gebeurt er wanneer $p$ onwaar is?

Beschouw:

\begin{center}
Als een geheel getal deelbaar is door $4$, dan is het even.
\end{center}

Neem $n=7$.

De voorwaarde

\[
7\text{ is deelbaar door }4
\]

is onwaar.

De uitspraak heeft niets beloofd over getallen die niet deelbaar zijn
door $4$.

Het getal $7$ kan de implicatie dus niet weerleggen.

Hetzelfde geldt voor $n=6$.

Ook daar is de voorwaarde ``deelbaar door $4$'' onwaar.
Dat $6$ wel even is, verandert daar niets aan.


\begin{opmerking}

Wanneer de voorwaarde $p$ onwaar is, kan de implicatie

\[
p\Rightarrow q
\]

niet worden weerlegd.

In de klassieke wiskundige logica beschouwen we de implicatie dan als
\textbf{waar}.

\end{opmerking}


Dit geeft de volledige waarheidstabel:

\[
\begin{array}{c|c|c}
p & q & p\Rightarrow q\\
\hline
\text{waar}   & \text{waar}   & \text{waar}\\
\text{waar}   & \text{onwaar} & \text{onwaar}\\
\text{onwaar} & \text{waar}   & \text{waar}\\
\text{onwaar} & \text{onwaar} & \text{waar}
\end{array}
\]


\begin{opmerking}

Een implicatie is slechts in één situatie onwaar:

\[
\important{
p\Rightarrow q
\text{ is onwaar precies wanneer }
p\text{ waar is en }q\text{ onwaar.}
}
\]

Dit is een van de belangrijkste regels uit dit hoofdstuk.

\end{opmerking}


\begin{oefening}[Waarheidswaarde van een implicatie]{\oefB\ESS}

Gegeven:

\[
p:\quad 10>5
\]

en

\[
q:\quad 10\text{ is oneven}.
\]

\begin{question}

Bepaal de waarheidswaarde van $p$.

\begin{oplossing}

$p$ is waar.

\end{oplossing}

\end{question}

\begin{question}

Bepaal de waarheidswaarde van $q$.

\begin{oplossing}

$q$ is onwaar.

\end{oplossing}

\end{question}

\begin{question}

Bepaal de waarheidswaarde van

\[
p\Rightarrow q.
\]

\begin{oplossing}

De voorwaarde $p$ is waar en de conclusie $q$ is onwaar.

Daarom is

\[
p\Rightarrow q
\]

onwaar.

\end{oplossing}

\end{question}

\end{oefening}


% ============================================================
\FVDsection{Een implicatie onderzoeken}
% ============================================================

Bij een uitspraak met veranderlijken moeten we goed onderscheiden tussen
een paar voorbeelden controleren en de algemene uitspraak onderzoeken.

Beschouw

\[
x>2\Rightarrow x^2>4
\]

met $x\in\mathbb{R}$.

Voor $x=3$ krijgen we

\[
3>2\Rightarrow 3^2>4.
\]

Beide delen zijn waar.

Ook voor $x=10$ werkt de uitspraak.

Maar enkele voorbeelden tonen nog niet aan dat de implicatie voor
\textbf{alle} toegelaten waarden geldt.

We moeten onderzoeken of er een waarde bestaat waarvoor de voorwaarde
waar is en de conclusie onwaar.


\begin{voorbeeld}[Een ware implicatie analyseren]

Onderzoek voor $x\in\mathbb{R}$:

\[
x>2\Rightarrow x^2>4.
\]

Als $x>2$, dan is $x$ positief.

Daarom mogen we de ongelijkheid $x>2$ kwadrateren zonder de orde te
veranderen:

\[
x^2>4.
\]

Telkens wanneer de voorwaarde waar is, is dus ook de conclusie waar.

De implicatie is waar.

\end{voorbeeld}


\begin{voorbeeld}[Een tegenvoorbeeld zoeken]

Onderzoek voor $x\in\mathbb{R}$:

\[
x^2>4\Rightarrow x>2.
\]

Op het eerste gezicht lijkt deze uitspraak aannemelijk.

Maar neem

\[
x=-3.
\]

Dan geldt

\[
(-3)^2=9>4,
\]

dus de voorwaarde is waar.

De conclusie

\[
-3>2
\]

is onwaar.

Daarom is de implicatie onwaar.

\end{voorbeeld}


\begin{opmerking}

Bij het onderzoeken van een implicatie is het vaak efficiënter om niet
willekeurig voorbeelden te proberen.

Vraag jezelf af:

\begin{center}
\textbf{Kan de voorwaarde waar zijn terwijl de conclusie onwaar is?}
\end{center}

Als dat kan, zoek dan gericht naar zo'n geval.

\end{opmerking}


\begin{oefening}[Een implicatie analyseren]{\oefU\ESS}

Beschouw voor $x\in\mathbb{R}$ de uitspraak

\[
x^2>9\Rightarrow x>3.
\]

\begin{question}

Onderzoek of de implicatie waar is.

\begin{oplossing}

De implicatie is onwaar.

We zoeken een waarde waarvoor

\[
x^2>9
\]

waar is maar

\[
x>3
\]

onwaar.

Neem bijvoorbeeld

\[
x=-4.
\]

Dan is

\[
(-4)^2=16>9,
\]

maar

\[
-4>3
\]

is onwaar.

Dus $x=-4$ is een tegenvoorbeeld.

\end{oplossing}

\end{question}

\begin{question}

Leg uit waarom het controleren van $x=4$, $x=5$ en $x=10$
niet volstaat om de algemene uitspraak te bewijzen.

\begin{oplossing}

Die waarden tonen alleen dat de implicatie voor enkele afzonderlijke
gevallen niet wordt geschonden.

Om een uitspraak over alle reële getallen te bewijzen, moeten alle
mogelijke waarden door een algemene redenering worden afgedekt.

Bovendien missen de gekozen voorbeelden de negatieve waarden die juist
tegenvoorbeelden opleveren.

\end{oplossing}

\end{question}

\end{oefening}


% ============================================================
\FVDsection{Voldoende en nodige voorwaarden}
% ============================================================

Een implicatie

\[
p\Rightarrow q
\]

kan op verschillende manieren in woorden worden uitgedrukt.

Beschouw:

\begin{center}
Als een geheel getal deelbaar is door $4$, dan is het even.
\end{center}

De eigenschap ``deelbaar zijn door $4$'' is voldoende om te besluiten
dat het getal even is.

We zeggen daarom:

\begin{center}
deelbaar zijn door $4$ is een \textbf{voldoende voorwaarde}
voor even zijn.
\end{center}

Omgekeerd: als een getal deelbaar is door $4$, dan moet het noodzakelijk
even zijn.

We zeggen daarom:

\begin{center}
even zijn is een \textbf{nodige voorwaarde}
voor deelbaar zijn door $4$.
\end{center}


\begin{definitie}

Als

\[
p\Rightarrow q,
\]

dan is

\begin{itemize}

  \item $p$ een \textbf{voldoende voorwaarde} voor $q$;

  \item $q$ een \textbf{nodige voorwaarde} voor $p$.

\end{itemize}

\end{definitie}


\begin{opmerking}

Een handige manier om dit te onthouden is:

\[
\boxed{
p\Rightarrow q
}
\]

betekent dat $p$ sterk genoeg is om $q$ te garanderen.

Dus $p$ is \textbf{voldoende}.

Tegelijk kan $p$ niet waar zijn zonder dat $q$ waar is.

Dus $q$ is \textbf{nodig}.

\end{opmerking}


\begin{voorbeeld}[Een meetkundig voorbeeld]

Beschouw:

\begin{center}
Als een vierhoek een vierkant is, dan is hij een rechthoek.
\end{center}

Dus:

\begin{itemize}

  \item vierkant zijn is een voldoende voorwaarde om rechthoek te zijn;

  \item rechthoek zijn is een nodige voorwaarde om een vierkant te zijn.

\end{itemize}

Rechthoek zijn is echter niet voldoende om een vierkant te zijn.

Een rechthoek met zijden $3$ en $5$ is bijvoorbeeld geen vierkant.

\end{voorbeeld}


\begin{oefening}[Nodig of voldoende?]{\oefU\ESS}

Gegeven is de ware implicatie:

\begin{center}
Als een geheel getal deelbaar is door $12$, dan is het deelbaar door $3$.
\end{center}

\begin{question}

Vul aan:

\begin{center}
Deelbaar zijn door $12$ is een \wordChoice{
\choice[correct]{voldoende}
\choice{nodige}
} voorwaarde voor deelbaar zijn door $3$.
\end{center}

\end{question}

\begin{question}

Vul aan:

\begin{center}
Deelbaar zijn door $3$ is een \wordChoice{
\choice{voldoende}
\choice[correct]{nodige}
} voorwaarde voor deelbaar zijn door $12$.
\end{center}

\end{question}

\begin{question}

Is deelbaar zijn door $3$ ook voldoende om deelbaar te zijn door $12$?
Verantwoord.

\begin{oplossing}

Nee.

Bijvoorbeeld $6$ is deelbaar door $3$, maar niet door $12$.

Dus deelbaar zijn door $3$ is niet voldoende om deelbaar te zijn door
$12$.

\end{oplossing}

\end{question}

\end{oefening}


% ============================================================
\FVDsection{De omgekeerde uitspraak}
% ============================================================

Beschouw opnieuw

\[
p\Rightarrow q.
\]

We kunnen de rollen van $p$ en $q$ verwisselen.

Dan krijgen we

\[
q\Rightarrow p.
\]

Dit noemen we de \textbf{omgekeerde uitspraak}.


\begin{definitie}

Bij de implicatie

\[
p\Rightarrow q
\]

is

\[
q\Rightarrow p
\]

de \textbf{omgekeerde uitspraak}.

\end{definitie}


\begin{voorbeeld}

De oorspronkelijke uitspraak is:

\begin{center}
Als een geheel getal deelbaar is door $4$, dan is het even.
\end{center}

De omgekeerde uitspraak is:

\begin{center}
Als een geheel getal even is, dan is het deelbaar door $4$.
\end{center}

De oorspronkelijke uitspraak is waar.

De omgekeerde uitspraak is onwaar.

Het getal

\[
6
\]

is immers even, maar niet deelbaar door $4$.

\end{voorbeeld}


\begin{opmerking}

Uit

\[
p\Rightarrow q
\]

mogen we in het algemeen \textbf{niet} besluiten dat

\[
q\Rightarrow p.
\]

Een implicatie en haar omgekeerde zijn verschillende uitspraken en
moeten afzonderlijk worden onderzocht.

\end{opmerking}


\begin{voorbeeld}[Een veelgemaakte denkfout]

We weten:

\begin{center}
Als een dier een hond is, dan is het een zoogdier.
\end{center}

Daaruit volgt niet:

\begin{center}
Als een dier een zoogdier is, dan is het een hond.
\end{center}

Een kat is een eenvoudig tegenvoorbeeld.

De fout ontstaat doordat de richting van de implicatie ongemerkt wordt
omgekeerd.

\end{voorbeeld}


\begin{oefening}[De richting doet ertoe]{\oefU\ESS}

Beschouw voor $x\in\mathbb{R}$:

\[
p:\quad x>3
\]

en

\[
q:\quad x^2>9.
\]

\begin{question}

Onderzoek

\[
p\Rightarrow q.
\]

\begin{oplossing}

De implicatie is waar.

Als $x>3$, dan is $x$ positief en volgt

\[
x^2>9.
\]

\end{oplossing}

\end{question}

\begin{question}

Formuleer de omgekeerde uitspraak.

\begin{oplossing}

\[
q\Rightarrow p,
\]

dus

\[
x^2>9\Rightarrow x>3.
\]

\end{oplossing}

\end{question}

\begin{question}

Onderzoek de omgekeerde uitspraak.

\begin{oplossing}

De omgekeerde uitspraak is onwaar.

Neem bijvoorbeeld

\[
x=-4.
\]

Dan is

\[
x^2=16>9,
\]

maar

\[
x>3
\]

is onwaar.

\end{oplossing}

\end{question}

\end{oefening}


% ============================================================
\FVDsection{De contrapositieve uitspraak}
% ============================================================

Naast de omgekeerde uitspraak is er nog een belangrijke uitspraak die
bij een implicatie hoort.

Vertrek van

\[
p\Rightarrow q.
\]

Ontken de conclusie en de voorwaarde en wissel hun plaats.

We krijgen

\[
\neg q\Rightarrow\neg p.
\]

Dit noemen we de \textbf{contrapositieve uitspraak}.


\begin{definitie}

Bij de implicatie

\[
p\Rightarrow q
\]

is

\[
\neg q\Rightarrow\neg p
\]

de \textbf{contrapositieve uitspraak}.

\end{definitie}


\begin{voorbeeld}

Beschouw:

\begin{center}
Als een geheel getal deelbaar is door $4$, dan is het even.
\end{center}

We nemen

\[
p:\quad n\text{ is deelbaar door }4
\]

en

\[
q:\quad n\text{ is even}.
\]

De contrapositieve uitspraak is

\[
\neg q\Rightarrow\neg p.
\]

In woorden:

\begin{center}
Als $n$ niet even is, dan is $n$ niet deelbaar door $4$.
\end{center}

Voor gehele getallen kunnen we ook schrijven:

\begin{center}
Als $n$ oneven is, dan is $n$ niet deelbaar door $4$.
\end{center}

\end{voorbeeld}


% ------------------------------------------------------------
\FVDsubsection{Implicatie en contrapositieve}
% ------------------------------------------------------------

Vergelijk de waarheidstabellen.

\[
\begin{array}{c|c|c|c|c}
p&q&p\Rightarrow q&\neg q&\neg q\Rightarrow\neg p\\
\hline
\text{waar}&\text{waar}&\text{waar}&\text{onwaar}&\text{waar}\\
\text{waar}&\text{onwaar}&\text{onwaar}&\text{waar}&\text{onwaar}\\
\text{onwaar}&\text{waar}&\text{waar}&\text{onwaar}&\text{waar}\\
\text{onwaar}&\text{onwaar}&\text{waar}&\text{waar}&\text{waar}
\end{array}
\]

De kolommen van

\[
p\Rightarrow q
\]

en

\[
\neg q\Rightarrow\neg p
\]

zijn identiek.

Daarom hebben beide uitspraken altijd dezelfde waarheidswaarde.


\begin{opmerking}

Een implicatie en haar contrapositieve zijn logisch gelijkwaardig:

\[
\important{
p\Rightarrow q
\quad\Longleftrightarrow\quad
\neg q\Rightarrow\neg p.
}
\]

Dit is fundamenteel voor een bewijstechniek die we later zullen leren:
\textbf{bewijs door contrapositie}.

\end{opmerking}


\begin{oefening}[Contrapositieve formuleren]{\oefB\ESS}

Beschouw:

\begin{center}
Als een geheel getal deelbaar is door $6$, dan is het deelbaar door $3$.
\end{center}

\begin{question}

Neem

\[
p:\quad n\text{ is deelbaar door }6
\]

en

\[
q:\quad n\text{ is deelbaar door }3.
\]

Welke formule stelt de contrapositieve voor?

\begin{multipleChoice}

  \choice{$q\Rightarrow p$}

  \choice{$\neg p\Rightarrow\neg q$}

  \choice[correct]{$\neg q\Rightarrow\neg p$}

  \choice{$p\Rightarrow\neg q$}

\end{multipleChoice}

\end{question}

\begin{question}

Formuleer de contrapositieve in woorden.

\begin{oplossing}

Als $n$ niet deelbaar is door $3$, dan is $n$ niet deelbaar door $6$.

\end{oplossing}

\end{question}

\end{oefening}


% ============================================================
\FVDsection{Equivalentie}
% ============================================================

Soms geldt een implicatie in beide richtingen.

Beschouw voor een geheel getal $n$:

\[
p:\quad n\text{ is even}
\]

en

\[
q:\quad n^2\text{ is even}.
\]

Dan geldt

\[
p\Rightarrow q.
\]

Als $n$ even is, dan is $n^2$ immers even.

Maar ook de omgekeerde uitspraak

\[
q\Rightarrow p
\]

is waar.

In zo'n geval zeggen we dat $p$ en $q$ equivalent zijn.


\begin{definitie}

De \textbf{equivalentie} van twee uitspraken $p$ en $q$ noteren we als

\[
p\Leftrightarrow q.
\]

We lezen dit als

\begin{center}
``$p$ als en slechts als $q$''
\end{center}

of

\begin{center}
``$p$ is equivalent met $q$''.
\end{center}

De equivalentie is waar wanneer $p$ en $q$ dezelfde waarheidswaarde
hebben.

\end{definitie}


De waarheidstabel is

\[
\begin{array}{c|c|c}
p&q&p\Leftrightarrow q\\
\hline
\text{waar}&\text{waar}&\text{waar}\\
\text{waar}&\text{onwaar}&\text{onwaar}\\
\text{onwaar}&\text{waar}&\text{onwaar}\\
\text{onwaar}&\text{onwaar}&\text{waar}
\end{array}
\]


\begin{opmerking}

Een equivalentie betekent dat beide implicaties gelden:

\[
\important{
p\Leftrightarrow q
\quad\text{betekent}\quad
(p\Rightarrow q)\land(q\Rightarrow p).
}
\]

We moeten dus twee richtingen controleren.

\end{opmerking}


% ------------------------------------------------------------
\FVDsubsection{Nodig én voldoende}
% ------------------------------------------------------------

Als

\[
p\Leftrightarrow q,
\]

dan geldt zowel

\[
p\Rightarrow q
\]

als

\[
q\Rightarrow p.
\]

Daarom is $p$ zowel een nodige als een voldoende voorwaarde voor $q$.

En omgekeerd is $q$ zowel nodig als voldoende voor $p$.


\begin{voorbeeld}

Voor $n\in\mathbb{Z}$ geldt:

\[
n\text{ is even}
\Leftrightarrow
n^2\text{ is even}.
\]

Dus:

\begin{center}
$n$ is even als en slechts als $n^2$ even is.
\end{center}

\end{voorbeeld}


\begin{voorbeeld}[Geen equivalentie]

Voor $n\in\mathbb{Z}$ geldt

\[
n\text{ is deelbaar door }4
\Rightarrow
n\text{ is even}.
\]

Maar de omgekeerde implicatie is niet altijd waar.

Daarom mogen we niet schrijven

\[
n\text{ is deelbaar door }4
\Leftrightarrow
n\text{ is even}.
\]

\end{voorbeeld}


\begin{oefening}[Implicatie of equivalentie?]{\oefU\ESS}

Onderzoek voor $n\in\mathbb{Z}$ de eigenschappen

\[
p:\quad n\text{ is deelbaar door }10
\]

en

\[
q:\quad n\text{ is deelbaar door }5.
\]

\begin{question}

Is

\[
p\Rightarrow q
\]

waar?

\begin{oplossing}

Ja.

Elk veelvoud van $10$ is ook een veelvoud van $5$.

\end{oplossing}

\end{question}

\begin{question}

Is

\[
q\Rightarrow p
\]

waar?

\begin{oplossing}

Nee.

Bijvoorbeeld $15$ is deelbaar door $5$, maar niet door $10$.

\end{oplossing}

\end{question}

\begin{question}

Mogen we schrijven

\[
p\Leftrightarrow q?
\]

\begin{oplossing}

Nee.

Voor een equivalentie moeten beide implicaties waar zijn.

Hier is alleen

\[
p\Rightarrow q
\]

waar.

\end{oplossing}

\end{question}

\end{oefening}


% ============================================================
\FVDsection{Verschillende formuleringen herkennen}
% ============================================================

In wiskundige teksten staat niet altijd letterlijk ``als \ldots, dan
\ldots''.

Dezelfde logische structuur kan op verschillende manieren worden
uitgedrukt.


\begin{voorbeeld}

De volgende uitspraken drukken dezelfde implicatie uit:

\begin{center}
Als $n$ deelbaar is door $4$, dan is $n$ even.
\end{center}

\begin{center}
$n$ deelbaar door $4$ is voldoende voor $n$ even.
\end{center}

\begin{center}
$n$ even is nodig voor $n$ deelbaar door $4$.
\end{center}

Symbolisch:

\[
p\Rightarrow q.
\]

\end{voorbeeld}


Voor een equivalentie kunnen we bijvoorbeeld schrijven:

\begin{center}
$p$ als en slechts als $q$,
\end{center}

\begin{center}
$p$ is nodig en voldoende voor $q$,
\end{center}

of

\[
p\Leftrightarrow q.
\]


\begin{oefening}[Logische taal interpreteren]{\oefU\ESS}

Beschouw de uitspraak:

\begin{center}
Deelbaar zijn door $6$ is voldoende om deelbaar te zijn door $3$.
\end{center}

Neem

\[
p:\quad n\text{ is deelbaar door }6
\]

en

\[
q:\quad n\text{ is deelbaar door }3.
\]

\begin{question}

Welke symbolische uitspraak hoort hierbij?

\begin{multipleChoice}

  \choice[correct]{$p\Rightarrow q$}

  \choice{$q\Rightarrow p$}

  \choice{$p\Leftrightarrow q$}

\end{multipleChoice}

\end{question}

\begin{question}

Welke van de volgende zinnen drukt dezelfde implicatie uit?

\begin{multipleChoice}

  \choice{Deelbaar zijn door $6$ is nodig voor deelbaar zijn door $3$.}

  \choice[correct]{Deelbaar zijn door $3$ is nodig voor deelbaar zijn door $6$.}

  \choice{Deelbaar zijn door $3$ is voldoende voor deelbaar zijn door $6$.}

\end{multipleChoice}

\end{question}

\end{oefening}


% ============================================================
\FVDsection{Logische structuur in toepassingen}
% ============================================================

Implicaties komen niet alleen in zuivere wiskunde voor.

Ook in informatica, elektronica en wetenschappen worden voorwaarden
voortdurend gebruikt.


\begin{voorbeeld}[Een temperatuuralarm]

Een regelsysteem bevat de volgende instructie:

\begin{center}
Als de temperatuur hoger is dan $80^\circ\mathrm{C}$,
dan wordt het alarm geactiveerd.
\end{center}

Neem

\[
p:\quad T>80^\circ\mathrm{C}
\]

en

\[
q:\quad \text{het alarm is actief}.
\]

De regel heeft de vorm

\[
p\Rightarrow q.
\]

Stel dat het alarm actief is.

Mogen we daaruit besluiten dat

\[
T>80^\circ\mathrm{C}?
\]

Niet noodzakelijk.

Het alarm kan misschien ook om een andere reden geactiveerd worden,
bijvoorbeeld door rookdetectie.

Uit

\[
p\Rightarrow q
\]

volgt immers niet automatisch

\[
q\Rightarrow p.
\]

\end{voorbeeld}


\begin{oefening}[Een fout besluit analyseren]{\oefU\ESS}

Een beveiligingssysteem werkt volgens de regel:

\begin{center}
Als de druk in een vat hoger is dan de veiligheidsgrens,
dan gaat het waarschuwingslampje branden.
\end{center}

Een leerling ziet dat het waarschuwingslampje brandt en besluit:

\begin{center}
``Dus de druk is hoger dan de veiligheidsgrens.''
\end{center}

\begin{question}

Is dit besluit logisch gerechtvaardigd op basis van de gegeven regel?

\begin{multipleChoice}

  \choice{Ja.}

  \choice[correct]{Nee.}

\end{multipleChoice}

\end{question}

\begin{question}

Leg uit welke logische fout wordt gemaakt.

\begin{oplossing}

Neem

\[
p:\quad \text{de druk is hoger dan de veiligheidsgrens}
\]

en

\[
q:\quad \text{het lampje brandt}.
\]

Gegeven is alleen

\[
p\Rightarrow q.
\]

De leerling gebruikt echter de omgekeerde uitspraak

\[
q\Rightarrow p.
\]

Die volgt niet automatisch uit de oorspronkelijke implicatie.

Het lampje zou bijvoorbeeld ook door een testfunctie of een andere
storing geactiveerd kunnen worden.

\end{oplossing}

\end{question}

\end{oefening}



\end{document}
